% ===========================================================================
% power.tex -- the unseen-drug arm's detection floor against the largest
% effect the mechanism channel could plausibly deliver.
% \input by Sections/05-Limitations.tex, in "It is underpowered by
% construction". Compiles against thesis_v2/preamble.tex and nothing else;
% no \includegraphics, no external file read at build time.
%
% WHY THIS FIGURE EXISTS. The passage it sits in makes its argument out of
% four numbers delivered in four separate sentences (0.052, 0.073, 0.079,
% 0.052-as-largest-effect). The conclusion -- the arm cannot detect the
% largest effect it could plausibly be measuring -- only follows once all
% four are on one scale. On one scale it is a picture: the bar ends before
% the threshold begins.
%
% ---------------------------------------------------------------------------
% EVERY NUMBER DRAWN, ITS SOURCE, AND THE THESIS SENTENCE IT MUST MATCH.
% Sibling power.json records the same list machine-readably.
%
%  threshold (dashed accent rule)      0.07897560995161473
%      RESULTS_cluster/CANONICAL_NUMBERS.json
%          .detection_limit.detectable_at_80pc_power
%      (same object carries .split = "unseen_drug", .n = 199,
%       .observed = 0.45889838436740604, .half_width = 0.05528292696613031;
%       the observed NIR and that half_width are NOT drawn -- this is a
%       difference axis and 0.4589 is an absolute NIR.)
%      Sections/05-Limitations.tex:310-311 "the smallest effect detectable
%      at $80\%$ power is $\approx 0.079$".                        MATCHES.
%
%  bar length                          0.05201514916327314
%      = 0.5520151491632731 - 0.500
%        RESULTS_cluster/CANONICAL_NUMBERS.json .means_all.moa_lookup
%        = RESULTS_cluster/re_v3.json .means.moa_lookup   (identical to 16 sf)
%        chance 0.500 is a definition: both frames are two-alternative
%        retrievals (preamble/frames.tex, Sections/03-Apparatus.tex).
%      Sections/05-Limitations.tex:311-314 "The largest effect the mechanism
%      channel could deliver is the distance from chance to
%      \texttt{moa\_lookup}, which scores $0.552$ pooled across splits
%      against a chance value of $0.500$, so $\approx 0.052$."      MATCHES.
%      0.552 and 0.500 are absolute NIR values and are therefore stated in
%      the caption only; the drawing carries the difference and nothing else.
%
%  point estimate (filled marker)     -0.03148175705606285
%      CANONICAL_NUMBERS.json .splits.unseen_drug.gap
%      = re_v3.json .means.generalization.unseen_drug.gap
%      Sections/05-Limitations.tex:322 \ci{-0.0315}{-0.1049}{+0.0419};
%      Sections/04c-reencoding.tex:791 table row, same estimate twice. MATCHES.
%
%  inner span (thick)  -0.08358077892636565 .. +0.020617264814239943
%      CANONICAL_NUMBERS.json .splits.unseen_drug.gap_ci_line
%      half-width = (0.020617264814239943 + 0.08358077892636565)/2
%                 = 0.052099021870302796  -> 0.052
%      Sections/05-Limitations.tex:308-310 "$\pm 0.052$ clustered two-way on
%      cell line and well". The key "gap_ci_line" is the JSON's own shorthand;
%      tab:generalisation's note says the interval is two-way cluster-robust
%      over cell lines and wells on min(G_line,G_well)-1 = 27 df, so the
%      figure uses the prose's two-way wording.                     MATCHES.
%
%  outer span (thin)   -0.10487258422550776 .. +0.04190907011338206
%      CANONICAL_NUMBERS.json .splits.unseen_drug.gap_ci_drug
%      half-width = (0.04190907011338206 + 0.10487258422550776)/2
%                 = 0.07339082716944491   -> 0.073
%      Sections/05-Limitations.tex:309-310 "$\pm 0.073$ clustered on the
%      drug"; the same interval is printed at 05:322 and at
%      04c-reencoding.tex:791.                                      MATCHES.
%
%  n = 199
%      CANONICAL_NUMBERS.json .splits.unseen_drug.n
%      Sections/04c-reencoding.tex:726 "$\num{199}$ unseen-drug" and the
%      tab:generalisation row "$199$ ($40$; $10$; $28$)".            MATCHES.
%
%  Both intervals are exactly symmetric about the point estimate: the
%  midpoint of gap_ci_line and the midpoint of gap_ci_drug each equal
%  -0.031481757056062853, i.e. .splits.unseen_drug.gap to 16 significant
%  figures. So "point estimate +- half-width" is the intervals themselves and
%  not a symmetrisation of them; nothing is being smoothed here.
%
% DELIBERATELY ABSENT, per the brief: any absolute NIR value inside the
% drawing (0.4589, 0.552, 0.500, 0.8365); any other split; any p-value; any
% legend (both spans are labelled where they end); any gridline.
%
% ---------------------------------------------------------------------------
% GEOMETRY. x is in NIR-difference units at x=383mm, so every coordinate in
% the source below is the datum itself and cannot drift from it by arithmetic
% done here. y is in mm.
%
%   axis          -0.12 .. +0.12   ->  -45.96mm .. +45.96mm   (91.92mm)
%   left column   44mm of flush-right label + 3mm gutter, east edge -48.96mm
%   total width   44 + 3 + 45.96 + 45.96 = 138.92mm, inside the 142.0mm
%                 measure (404.03pt, preamble.tex:793). Measured on the
%                 harness build: see _harness_power.log "HARNESS tikz width".
%   bar tip       0.05201515 * 383 = 19.92mm
%   threshold     0.07897561 * 383 = 30.25mm
%   -> the bar stops 10.33mm short of the rule. That clearance is the whole
%      argument, so it is stated here as a measured length and not judged by
%      eye; at 142mm it is about 7 per cent of the measure and reads at once.
%
%   rows, y in mm   bar      33.5 (rectangle 30.7..36.3)
%                   estimate 13.5 (inner caps 11.1..15.9, outer 12.0..15.0)
%                   axis      0.0, ticks to -1.4, axis title at -7.6
%   detection rule  -1.4 .. 42.5, its label sitting on top of it at 42.6
%   zero rule       -1.4 .. 16.4 -- see note at the draw itself
%
%   MEASURED on the harness build (_harness_power.log):
%     HARNESS textwidth  = 404.02913pt = 142.00mm
%     HARNESS tikz width = 395.70499pt = 139.07mm   <- 2.93mm inside the
%     HARNESS tikz height= 173.09203pt =  60.83mm      measure, 0 overfull
%   and in the whole-thesis build the only Overfull \hbox in main.log is the
%   pre-existing 1.35pt one in figs/fig-frame.tex, not this file.
%
% COLOUR. accent (#0B5394) is spent on exactly two marks, the bar and the
% detection floor, because those two are the comparison the passage makes.
% The observed estimate is ink, the axis and the zero rule rulegrey, all
% descriptive text quietink and every numeral ink. Nothing is told apart by
% colour that is not also told apart by position, weight or shape: the two
% nested spans differ in thickness and in length, not in hue, and the figure
% survives greyscale printing.
% ===========================================================================
\begin{figure}[htbp]
% MARGINALIA (06-figures.md): wider than the 124mm text column, so it
% sits in fullwidth; the caption below spans the text column only.
\begin{fullwidth}\raggedright
\begin{tikzpicture}[
    x=383mm, y=1mm, font=\footnotesize,
    rowlbl/.style  ={text=quietink, align=flush right, text width=44mm,
                     inner sep=0pt, anchor=east, xshift=-3mm},
    ticklbl/.style ={font=\scriptsize, text=quietink, inner sep=0pt,
                     anchor=north},
    axtitle/.style ={font=\scriptsize, text=quietink, inner sep=0pt,
                     anchor=north},
    note/.style    ={text=quietink, inner sep=0pt}]

  % ---- 1. the difference axis ---------------------------------------------
  % Zero, not 0.50, is the reference: this axis carries differences, so the
  % chance level lives inside each quantity's definition and not on the rule.
  \draw[rulegrey, line width=0.5pt] (-0.12,0) -- (0.12,0);
  \foreach \t/\lab in {%
      -0.10/$-0.10$, -0.05/$-0.05$, 0/$0$, 0.05/$+0.05$, 0.10/$+0.10$}{%
    \draw[rulegrey, line width=0.5pt] (\t,0) -- (\t,-1.4);
    \node[ticklbl] at (\t,-2.1) {\lab};}
  \node[axtitle] at (0,-7.6) {difference in $\NIR$\enspace($0$ = no difference)};

  % ---- 2. zero -------------------------------------------------------------
  % The rule stops at 16.4, just under the inner span's label. It does not
  % need to reach the bar row: the bar's own left edge IS x=0 and carries the
  % position up on its own. Running it higher would have to be struck through
  % the "+-0.052" label, which is 46mm wide and must end flush with the inner
  % span's right cap at +0.0206, so it cannot avoid crossing zero. A knocked-
  % out rule reads as a broken line; a shorter one reads as a rule.
  \draw[rulegrey, line width=0.4pt] (0,-1.4) -- (0,16.4);

  % ---- 3. the detection floor ---------------------------------------------
  \draw[accent, line width=0.9pt, dash pattern=on 2.2pt off 1.6pt]
    (0.07897561,-1.4) -- (0.07897561,42.5);
  \node[note, anchor=south east, align=flush right, xshift=-1.0mm]
    at (0.07897561,42.6)
    {smallest effect this arm can\\detect at $80\%$ power,
     \textcolor{ink}{$0.079$}};

  % ---- 4. the largest effect the mechanism channel could deliver ----------
  % It starts at zero because it IS a distance from a null, and it ends
  % where it ends. The dotted arrow is the shortfall: it is drawn without a
  % number on it because the thesis states 0.052 and 0.079 and does not
  % state their difference.
  \fill[accent!20]                (0,30.7) rectangle (0.05201515,36.3);
  \draw[accent, line width=0.7pt] (0,30.7) rectangle (0.05201515,36.3);
  \draw[quietink, line width=0.6pt, dotted, line cap=round,
        -{Latex[length=1.5mm,width=1.3mm]}]
    ([xshift=1.4mm]0.05201515,33.5) -- ([xshift=-0.8mm]0.07897561,33.5);
  \node[rowlbl] at (-0.12,33.5)
    {largest effect the\\mechanism channel\\could deliver\\[0.35ex]
     \textcolor{ink}{$0.052$}};

  % ---- 5. the arm's own estimate, with both cluster keys ------------------
  % Two nested spans, same point estimate, different unit of generalisation.
  % Thin/long = clustered on drug; thick/short = clustered two-way on cell
  % line and well. Told apart by weight and by extent, never by colour.
  % Each label ends FLUSH with the span end it names -- the inner one at its
  % right cap, above; the outer one at its left cap, below -- so the two are
  % separated diagonally and neither needs a legend or a leader line.
  \draw[ink, line width=0.5pt] (-0.10487258,13.5) -- (0.04190907,13.5);
  \draw[ink, line width=0.5pt] (-0.10487258,12.0) -- (-0.10487258,15.0);
  \draw[ink, line width=0.5pt] ( 0.04190907,12.0) -- ( 0.04190907,15.0);
  \draw[ink, line width=1.3pt] (-0.08358078,13.5) -- (0.02061726,13.5);
  \draw[ink, line width=1.0pt] (-0.08358078,11.1) -- (-0.08358078,15.9);
  \draw[ink, line width=1.0pt] ( 0.02061726,11.1) -- ( 0.02061726,15.9);
  \fill[white] (-0.03148176,13.5) circle (2.1pt);
  \fill[ink]   (-0.03148176,13.5) circle (1.6pt);

  \node[note, anchor=south east, align=flush right] at (0.02061726,16.6)
    {$\pm 0.052$ clustered two-way\\on cell line and well};
  % broken over two lines, not one, so that its right edge stops at -0.026
  % and the label clears the zero rule instead of being knocked out of it.
  \node[note, anchor=north west, align=flush left] at (-0.10487258,10.4)
    {$\pm 0.073$ clustered\\on drug};

  \node[rowlbl] at (-0.12,13.5)
    {observed \texttt{unseen\_drug} gap\\($n = 199$)\\[0.35ex]
     \textcolor{ink}{$-0.0315$}};
\end{tikzpicture}
\caption[The detection floor sits above the largest plausible
effect]{\captiontitle{The bar ends before the threshold begins.} Every quantity
here is a difference in $\NIR$ on one scale; no absolute $\NIR$ is plotted.
The dashed rule is the smallest effect the \texttt{unseen\_drug} arm can
detect at $80\%$ power, $0.079$, on the $n = \num{199}$ conditions that arm
scores. The bar is the largest effect the mechanism channel could deliver ---
the distance from chance to the mechanism-only lookup, which scores $0.552$
pooled across splits against a chance value of $0.500$, so $\approx 0.052$ ---
and it stops short of the rule. Beneath it is the arm's own estimate: model
$\NIR$ minus the geometry-defined \texttt{orth} scramble comparator on the
same conditions, $-0.0315$, drawn with the thick span the $\pm 0.052$
half-width clustered two-way on cell line and well and the thin span the
$\pm 0.073$ half-width clustered on the drug. Both cover zero, and the
half-width on the tighter of the two is itself the size of the bar. The arm
is licensed as an instrument control and is not evidence that unseen-drug
transfer is zero: an effect of the largest size the mechanism channel could
deliver would sit inside that same interval, so what it supplies is the
absence of a contradiction rather than a confirmed null.}
\label{fig:power}
\end{fullwidth}
\end{figure}
